% TOP_PROBLEM: {"id": 5100056, "title": "Elliptic-billiard invariant k_{813}", "category_id": 11, "category": {"id": 11, "name": "dynamical_systems", "display_name": "Dynamical Systems", "description": "Problems about long-term behavior of deterministic systems, Hamiltonian dynamics, and periodic orbits.", "slug": "dynamical-systems", "order_index": 11, "created_at": "2026-07-31T15:26:25.671Z"}, "status": "open"}
% FIRST_POSTED: null
% ENTRY_KIND: statement_only
% Imported from: batch02.tex; UnsolvedMath ID 5100056
\documentclass[11pt]{article}
\usepackage{amsmath,amssymb,mathtools}
\usepackage{fontspec,unicode-math}
\setmainfont{Latin Modern Roman}
\setsansfont{Latin Modern Sans}
\setmathfont{Latin Modern Math}
\usepackage{xurl,hyperref}
\newcommand{\sref}[2]{\hyperlink{source:#1:#2}{[S#2]}}
\newcommand{\cataloguescope}[1]{\paragraph{Mathematical question.}#1}
\begin{document}
% UnsolvedMath ID 5100056; source catalogue code AMR-050-0056
\setcounter{section}{5100055}
\section{Elliptic-billiard invariant $k_{813}$}\label{5100056}
{\small\sffamily UnsolvedMath ID 5100056 · Dynamical Systems}
\subsection{Definitions and mathematical statement}

\paragraph{Shared notation.}$\mathbb N=\{1,2,\ldots\}$, and $\mathbb Z,\mathbb Q,\mathbb R,\mathbb C$ denote the integers, rationals, reals and complex numbers. Any other conventions are specified locally.


Fix $a>b>0$ and the ellipse $E:x^2/a^2+y^2/b^2=1$, with foci $f_1,f_2$. Fix an inner confocal elliptic caustic $C$ admitting a Poncelet family of $N$-periodic billiard polygons $P=(P_i)_{i=1}^N$: vertices lie on $E$, sides are tangent to $C$, and consecutive sides satisfy the billiard reflection law. Indices are cyclic and polygons retain their cyclic order. The family and its caustic stay fixed while the starting vertex varies. All areas are signed: for $W=(W_i)$,
\[\mathcal A(W)=\frac12\sum_{i=1}^N(x_i y_{i+1}-x_{i+1}y_i),\qquad W_i=(x_i,y_i).\]
Write $A=\mathcal A(P)$. Unit-circle inversion about $F$ is $I_F(X)=F+(X-F)/\lVert X-F\rVert^2$; set $P_j^\dagger=(I_{f_j}(P_i))$ and $A_j^\dagger=\mathcal A(P_j^\dagger)$, for $j=1,2$. Quotients are considered where their denominators and the polygon constructions are defined. For a polygon $W=(W_i)$ and point $F$, its antipedal polygon has as vertices the consecutive intersections of the lines through $W_i$ perpendicular to $W_i-F$. Write $A_{j,\mathrm{ant}}$ for the signed area of the antipedal polygon of $P$ with respect to $f_j$. The polar polygon with respect to $f_j$ is the antipedal of $P_j^\dagger$ with respect to $f_j$. Its signed area is $A_{j,\mathrm{pol}}$ and its internal angles are $\psi_{j,i}$. This polygon has a fixed circular caustic; let $O_{p,j}$ be its center. The dual polygon is the vertexwise inversion of the polar polygon about $O_{p,j}$, using a fixed unit radius. Denote its signed area by $A_{j,\mathrm{dual}}$ and its side lengths by $w_{j,i}$ (the source writes $w_i$ for a chosen focus). The notation $A_{j,\mathrm{inv}}$ in Table 9 denotes the area $A_j^\dagger$ of the inversive polygon.
\paragraph{Statement recorded in the supplied source.}
Prove that, for all $N$, the quantity $A_{j,\mathrm{pol}}/A_j^\dagger$ is constant as $P$ varies through the family. \sref{5100056}{2}
\subsection{Short English statement}
Show that, for all $N$, the quantity $A_{j,\mathrm{pol}}/A_j^\dagger$ is constant as $P$ varies through the family. The constant may depend on the fixed ellipse and caustic. These are the experimentally recorded assertions of the cited 2020 version.
\subsection{Sources}
\begin{description}
\item[\hypertarget{source:5100056:1}{S1}] ULAM.AI and the UnsolvedMath Contributors, UnsolvedMath: A Curated Collection of Open Mathematics Problems, record 5100056 (AMR-050-0056). CC BY 4.0. \url{https://huggingface.co/datasets/ulamai/UnsolvedMath}.
\item[\hypertarget{source:5100056:2}{S2}] Dan Reznik, Ronaldo Garcia and Jair Koiller, \emph{Eighty New Invariants in the Elliptic Billiard}, arXiv:2004.12497v11 (29 October 2020), Section 2 and Section 3.9, Table 9, row $k_{813}$, pp. 3, 9--12. \url{https://arxiv.org/abs/2004.12497v11}.
\end{description}
\label{end:5100056}
\end{document}
